\chapter{FAIR Records, CWL Export, Operations and the Capstone}
\label{ch:fair-cwl}

A workflow that produces a correct result has done half its work. The other
half is to let someone else find the result, trust it and run it again. This
chapter covers the last four lectures of the course. It shows FAIR metadata
(Lecture~45), export to the Common Workflow Language (Lecture~46), the
operational helpers that keep a deployment honest (Lecture~47), and the
capstone that puts every task type into one graph (Lecture~48).

\section{FAIR Metadata}

FAIR metadata is opt-in. A workflow that does not call \texttt{enable\_fair}
writes nothing extra. With a \texttt{FairProfile} (title, description,
creators, licence, keywords, access policy) and outputs declared as
\texttt{Artifact}s, every run writes a bundle under
\texttt{.dagon/fair/<workflow>/<run>/}. The bundle holds nine files: an
RO-Crate, a W3C PROV document, DataCite and CodeMeta records, the run record,
a fairness report checked against F1, F2, A1, I1, I2 and R1.1, a SHA-256
manifest, and Markdown and HTML reports.

\lstinputlisting[caption={Lecture 45: a FAIR run (\texttt{examples/45-fair/fair.py}).},label={lst:fair}]{../../examples/45-fair/fair.py}

\lstinputlisting[style=bash,caption={Recorded output of Lecture 45.}]{../../examples/45-fair/output/fair.txt}

\begin{cautionbox}[Credentials are never exported]
The configuration of Lecture~45 contains a (fake) API key. Values whose keys
match the redaction patterns (\texttt{TOKEN}, \texttt{SECRET},
\texttt{PASSWORD}, \texttt{PRIVATE\_KEY} and the cloud-provider prefixes) are
removed before anything is written, and the example checks that the key
appears in no file of the bundle.
\end{cautionbox}

\section{CWL Export}

\texttt{saveAsCWL} writes a CWL v1.2 workflow. File references become step
inputs. An ordering edge with no file, such as \texttt{notify} after
\texttt{mean\_ISC}, becomes a boolean \texttt{completed} output of the
producer and an \texttt{after\_<producer>} input of the consumer. The export
is deterministic: two exports of the same graph are byte-identical, and no
scratch path appears in them. Container tasks gain a
\texttt{DockerRequirement} hint.

\lstinputlisting[caption={Lecture 46: export and validation (\texttt{examples/46-cwl/cwl\_export.py}).},label={lst:cwl}]{../../examples/46-cwl/cwl_export.py}

\lstinputlisting[style=bash,caption={Recorded output of Lecture 46.}]{../../examples/46-cwl/output/cwl_export.txt}

\section{Operating DAGonStar}

Lecture~47 collects the helpers an operator relies on. With
\texttt{[dagon\_service]} enabled, a workflow checks the service, registers
itself and registers each task. The JSON form keeps explicit edges
(\Cref{ch:course-pipeline}). The runner environment keeps the user's
\texttt{PYTHONPATH} entries first and appends the package root. The Bash used
for launchers can be overridden with \texttt{DAGON\_BASH}. The packaged
\texttt{dagon-agent-runner} is installed as a console script.

\lstinputlisting[style=bash,caption={Recorded output of Lecture 47 (\texttt{examples/47-operations/operations.py}).}]{../../examples/47-operations/output/operations.txt}

\section{The Capstone}

The capstone builds one graph of eighteen tasks that uses all sixteen
registered task types: checkpoint, batch, Slurm, cloud, Docker, Kubernetes,
Apptainer, Nomad, LLM, native, web, FaaS, IoT, agentic, meta and volume. It runs under
portable emulation with checkpointing and FAIR recording enabled. Every
network-facing adapter is replaced by a function that fails the test if it is
called. The run must therefore succeed without a single live call, and the
CWL export must still describe all eighteen steps.

\lstinputlisting[caption={Lecture 48: the capstone (\texttt{examples/48-capstone/capstone.py}).},label={lst:capstone}]{../../examples/48-capstone/capstone.py}

\lstinputlisting[style=bash,caption={Recorded output of Lecture 48.}]{../../examples/48-capstone/output/capstone.txt}

\begin{practicebox}
The capstone is also a template for continuous integration. A workflow that
runs end to end under emulation, with every live adapter guarded, catches
broken references, bad configurations and contract errors on every commit.
It needs no cluster, no credential and no quota.
\end{practicebox}

\section{Chapter Summary}

\begin{itemize}[leftmargin=1.4em]
  \item FAIR bundles describe declared outputs, licences, provenance and
    checksums, and never contain credentials.
  \item CWL exports are deterministic, path-free and keep ordering edges.
  \item Operational helpers make service registration, environments and
    shells explicit.
  \item The capstone proves that the whole task-type catalogue composes into
    one checkpointed, recorded and exportable graph.
\end{itemize}

\begin{exercise}
Enable FAIR recording on the station pipeline of \Cref{ch:course-pipeline},
declare the report as an artifact under CC-BY-4.0, and read the fairness
report. Which checks pass, and which require information that the workflow
does not have?
\end{exercise}
